


# This is text file that contains all of the shell comands ran on EMR shell that were used to complete part 3 and some of 4 for the project





ssh -i /Users/maxmarte/Downloads/emr-lab-key.pem hadoop@10.99.17.21


hdfs dfs -mkdir -p /user/hadoop/bank_staging

#loading staging tables to hdfs from S3 where they have to already be located
hdfs dfs -cp s3://jpmorgan2025/account_table.csv      /user/hadoop/bank_staging/
hdfs dfs -cp s3://jpmorgan2025/customer_table.csv     /user/hadoop/bank_staging/
hdfs dfs -cp s3://jpmorgan2025/loan_table.csv         /user/hadoop/bank_staging/
hdfs dfs -cp s3://jpmorgan2025/transaction_table.csv  /user/hadoop/bank_staging/


Hive

CREATE DATABASE IF NOT EXISTS bank_dw;
USE bank_dw;

Create staging tables

CREATE TABLE IF NOT EXISTS stg_customer (
  customer_id         STRING,
  age                 INT,
  job                 STRING,
  marital             STRING,
  education           STRING,
  demographic_balance DOUBLE,
  housing_loan        INT,
  personal_loan       INT,
  cc_balance          DOUBLE,
  credit_limit        DOUBLE,
  purchases           DOUBLE,
  payments            DOUBLE
)
ROW FORMAT SERDE 'org.apache.hadoop.hive.serde2.OpenCSVSerde'
WITH SERDEPROPERTIES (
  "separatorChar" = ";",
  "quoteChar"     = "\""
)
STORED AS TEXTFILE;

LOAD DATA INPATH '/user/hadoop/bank_staging/customer_table.csv'
INTO TABLE stg_customer;

CREATE TABLE IF NOT EXISTS stg_account (
  account_id   STRING,
  customer_id  STRING,
  account_type STRING,
  open_date    STRING
)
ROW FORMAT SERDE 'org.apache.hadoop.hive.serde2.OpenCSVSerde'
WITH SERDEPROPERTIES (
  "separatorChar" = ";",
  "quoteChar"     = "\""
)
STORED AS TEXTFILE;

LOAD DATA INPATH '/user/hadoop/bank_staging/account_table.csv'
INTO TABLE stg_account;

CREATE TABLE IF NOT EXISTS stg_loan (
  loan_id     STRING,
  customer_id STRING,
  loan_type   STRING,
  loan_amount DOUBLE,
  start_date  STRING,
  status      STRING
)
ROW FORMAT SERDE 'org.apache.hadoop.hive.serde2.OpenCSVSerde'
WITH SERDEPROPERTIES (
  "separatorChar" = ";",
  "quoteChar"     = "\""
)
STORED AS TEXTFILE;

LOAD DATA INPATH '/user/hadoop/bank_staging/loan_table.csv'
INTO TABLE stg_loan;

CREATE TABLE IF NOT EXISTS stg_transaction (
  transaction_id INT,
  customer_id    STRING,
  account_id     STRING,
  transaction_ts STRING,
  amount         DOUBLE,
  category       STRING,
  is_fraud       INT
)
ROW FORMAT SERDE 'org.apache.hadoop.hive.serde2.OpenCSVSerde'
WITH SERDEPROPERTIES (
  "separatorChar" = ";",
  "quoteChar"     = "\""
)
STORED AS TEXTFILE;

LOAD DATA INPATH '/user/hadoop/bank_staging/transaction_table.csv'
INTO TABLE stg_transaction;

#ETL Extract + Load into staging is done


#Now onto the transform part

CREATE TABLE IF NOT EXISTS dim_customer (
  customer_id         STRING,
  age                 INT,
  job                 STRING,
  marital             STRING,
  education           STRING,
  demographic_balance DOUBLE,
  housing_loan        INT,
  personal_loan       INT,
  cc_balance          DOUBLE,
  credit_limit        DOUBLE,
  purchases           DOUBLE,
  payments            DOUBLE
)
STORED AS TEXTFILE;

INSERT OVERWRITE TABLE dim_customer
SELECT
  customer_id,
  age,
  job,
  marital,
  education,
  demographic_balance,
  housing_loan,
  personal_loan,
  cc_balance,
  credit_limit,
  purchases,
  payments
FROM stg_customer;



CREATE TABLE IF NOT EXISTS dim_account (
  account_id   STRING,
  customer_id  STRING,
  account_type STRING,
  open_date    STRING      -- keep as string for simplicity
)
STORED AS TEXTFILE;

INSERT OVERWRITE TABLE dim_account
SELECT
  account_id,
  customer_id,
  account_type,
  open_date
FROM stg_account;






CREATE TABLE IF NOT EXISTS dim_category (
  category STRING
)
STORED AS TEXTFILE;

INSERT OVERWRITE TABLE dim_category
SELECT DISTINCT category
FROM stg_transaction;





CREATE TABLE IF NOT EXISTS dim_date (
  date_id STRING,   -- 'yyyy-MM-dd'
  year    INT,
  month   INT,
  day     INT
)
STORED AS TEXTFILE;

INSERT OVERWRITE TABLE dim_date
SELECT
  date_format(parsed_ts, 'yyyy-MM-dd') AS date_id,
  year(parsed_ts)                       AS year,
  month(parsed_ts)                      AS month,
  day(parsed_ts)                        AS day
FROM (
  SELECT
    from_unixtime(
      unix_timestamp(transaction_ts,'yyyy-MM-dd HH:mm:ss')
    ) AS parsed_ts
  FROM stg_transaction
) t
GROUP BY
  date_format(parsed_ts, 'yyyy-MM-dd'),
  year(parsed_ts),
  month(parsed_ts),
  day(parsed_ts);



SET hive.exec.dynamic.partition = true;
SET hive.exec.dynamic.partition.mode = nonstrict;

INSERT OVERWRITE TABLE fact_transaction
PARTITION (year, month)
SELECT
  t.transaction_id,
  t.customer_id,
  t.account_id,
  date_format(parsed_ts, 'yyyy-MM-dd') AS date_id,
  t.amount,
  t.category,
  t.is_fraud,
  year(parsed_ts)                      AS year,
  month(parsed_ts)                     AS month
FROM (
  SELECT
    transaction_id,
    customer_id,
    account_id,
    from_unixtime(
      unix_timestamp(transaction_ts,'yyyy-MM-dd HH:mm:ss')
    ) AS parsed_ts,
    amount,
    category,
    is_fraud
  FROM stg_transaction
) t;

#Make big table for queries

USE bank_dw;

CREATE TABLE customer_spend_analytics AS
SELECT
    ft.customer_id,
    dc.age,
    dc.job,
    dc.marital,
    dc.education,
    ft.category,
    dd.year,
    dd.month,
    ft.amount
FROM fact_transaction ft
JOIN dim_customer dc 
    ON ft.customer_id = dc.customer_id
JOIN dim_date dd
    ON ft.date_id = dd.date_id;


#This is the python file that is written out to hive that the front end will call when a user inputs an age range. It will query based on that age range
# and with the results of those queries will create visualizations in python and then put on the front end


nano spending_analyzer_2.py

#!/usr/bin/env python3
"""
Customer Spending Behavior Analyzer (No Date Filter)
Pretty visuals, NO Monthly Spending Trend subplot.

Generates ONE PNG:
    /home/hadoop/spending_analysis_<min_age>_<max_age>.png

Panels:
  [1] Total Spending by Purchase Type
  [2] Average Spending by Job Type (Top 10)
  [3] Average Spending by Education Level
  [4] Average Spending by Marital Status
  [5] Transaction Distribution by Category (pie)

Monthly spending trend is still computed for the text report, but not plotted.
"""

from pyhive import hive
import pandas as pd
import matplotlib
matplotlib.use('Agg')  # Use non-interactive backend
import matplotlib.pyplot as plt
import seaborn as sns
from matplotlib.ticker import FuncFormatter
from datetime import datetime
import warnings
warnings.filterwarnings('ignore')


class SpendingAnalyzer:
    def __init__(self, host='localhost', port=10000, database='default'):
        self.host = host
        self.port = port
        self.database = database
        self.connection = None

    def connect(self):
        """Establish connection to Hive"""
        try:
            self.connection = hive.Connection(
                host=self.host,
                port=self.port,
                database=self.database,
                auth='NONE'
            )
            print("✓ Connected to Hive successfully")
            return True
        except Exception as e:
            print(f"✗ Connection failed: {e}")
            return False

    def get_age_range_data(self, min_age, max_age):
        """Retrieve ALL spending data for age range (no date filter)."""
        query = f"""
        SELECT 
            customer_id,
            age,
            job,
            marital,
            education,
            category,
            year,
            month,
            amount
        FROM customer_spend_analytics
        WHERE age BETWEEN {min_age} AND {max_age}
        ORDER BY year, month
        """

        print(f"\n🔍 Querying Hive for data...")
        print(f"   Age range: {min_age}-{max_age}")

        try:
            df = pd.read_sql(query, self.connection)
            if len(df) > 0:
                years = df['year'].unique()
                print(f"✓ Retrieved {len(df)} transactions "
                      f"from {df['customer_id'].nunique()} unique customers")
                print(f"   Data spans: {years.min()}-{years.max()}")
            return df
        except Exception as e:
            print(f"✗ Query failed: {e}")
            return None

    def analyze_spending_behavior(self, df):
        """Compute summary tables used for visuals + report."""
        print("\n📊 Analyzing spending behavior...")

        analysis = {}

        # Category analysis
        analysis['avg_by_category'] = df.groupby('category')['amount'].agg(
            ['mean', 'median', 'sum', 'count']
        )
        analysis['avg_by_category'].columns = ['Average', 'Median', 'Total', 'Count']

        # Job type
        analysis['avg_by_job'] = df.groupby('job')['amount'].agg(
            ['mean', 'median', 'sum', 'count']
        )
        analysis['avg_by_job'].columns = ['Average', 'Median', 'Total', 'Count']
        analysis['avg_by_job'] = analysis['avg_by_job'].sort_values('Average', ascending=False)

        # Education level
        analysis['avg_by_education'] = df.groupby('education')['amount'].agg(
            ['mean', 'median', 'sum', 'count']
        )
        analysis['avg_by_education'].columns = ['Average', 'Median', 'Total', 'Count']
        analysis['avg_by_education'] = analysis['avg_by_education'].sort_values('Average', ascending=False)

        # Marital status
        analysis['avg_by_marital'] = df.groupby('marital')['amount'].agg(
            ['mean', 'median', 'sum', 'count']
        )
        analysis['avg_by_marital'].columns = ['Average', 'Median', 'Total', 'Count']

        # Monthly trend (for text only)
        df['year_month'] = df['year'].astype(str) + '-' + df['month'].astype(str).str.zfill(2)
        analysis['monthly_trend'] = df.groupby('year_month')['amount'].agg(
            ['sum', 'mean', 'count']
        )
        analysis['monthly_trend'].columns = ['Total', 'Average', 'Transaction_Count']

        # Category trend over time (for possible future use)
        analysis['category_trend'] = df.groupby(['year_month', 'category'])['amount'] \
                                       .sum().unstack(fill_value=0)

        # Overall stats
        analysis['overall_stats'] = {
            'total_spending': df['amount'].sum(),
            'avg_transaction': df['amount'].mean(),
            'median_transaction': df['amount'].median(),
            'total_transactions': len(df),
            'unique_customers': df['customer_id'].nunique(),
            'avg_spending_per_customer': df.groupby('customer_id')['amount'].sum().mean()
        }

        print("✓ Analysis complete")
        return analysis

    def generate_text_report(self, age_range, analysis):
        """Print the text report to stdout."""
        min_age, max_age = age_range

        print("\n" + "="*80)
        print("CUSTOMER SPENDING BEHAVIOR ANALYSIS REPORT")
        print(f"Age Range: {min_age}-{max_age} years old")
        print(f"Generated: {datetime.now().strftime('%Y-%m-%d %H:%M:%S')}")
        print("="*80)

        stats = analysis['overall_stats']
        print("\n📈 OVERALL STATISTICS")
        print("-" * 80)
        print(f"Total Spending:              ${stats['total_spending']:,.2f}")
        print(f"Average Transaction:         ${stats['avg_transaction']:,.2f}")
        print(f"Median Transaction:          ${stats['median_transaction']:,.2f}")
        print(f"Total Transactions:          {stats['total_transactions']:,}")
        print(f"Unique Customers:            {stats['unique_customers']:,}")
        print(f"Avg Spending per Customer:   ${stats['avg_spending_per_customer']:,.2f}")

        print("\n💳 SPENDING BY PURCHASE TYPE")
        print("-" * 80)
        print(analysis['avg_by_category'].to_string())

        print("\n👔 SPENDING BY JOB TYPE (Top 10)")
        print("-" * 80)
        print(analysis['avg_by_job'].head(10).to_string())

        print("\n🎓 SPENDING BY EDUCATION LEVEL")
        print("-" * 80)
        print(analysis['avg_by_education'].to_string())

        print("\n💑 SPENDING BY MARITAL STATUS")
        print("-" * 80)
        print(analysis['avg_by_marital'].to_string())

        print("\n📅 MONTHLY SPENDING TREND")
        print("-" * 80)
        print(analysis['monthly_trend'].to_string())

        print("\n💡 KEY INSIGHTS")
        print("-" * 80)

        highest_category = analysis['avg_by_category']['Total'].idxmax()
        highest_amount = analysis['avg_by_category']['Total'].max()
        print(f"• Highest total spending category: {highest_category} "
              f"(${highest_amount:,.2f})")

        most_common = analysis['avg_by_category']['Count'].idxmax()
        most_common_count = analysis['avg_by_category']['Count'].max()
        print(f"• Most common purchase type: {most_common} "
              f"({most_common_count:,} transactions)")

        top_job = analysis['avg_by_job'].index[0]
        top_job_avg = analysis['avg_by_job']['Average'].iloc[0]
        print(f"• Highest average spending job: {top_job} "
              f"(${top_job_avg:,.2f} per transaction)")

        trend_data = analysis['monthly_trend']['Total']
        if len(trend_data) >= 6:
            recent_avg = trend_data.tail(3).mean()
            early_avg = trend_data.head(3).mean()
            change_pct = ((recent_avg - early_avg) / early_avg) * 100
            trend_direction = "increased" if change_pct > 0 else "decreased"
            print(
                f"• Spending trend: {trend_direction} by {abs(change_pct):.1f}% "
                f"(comparing recent vs earlier months)"
            )

        print("\n" + "="*80)

    def create_visualizations(self, age_range, analysis, output_dir='/home/hadoop/'):
        """
        Create and save ONE summary PNG.
        NO Monthly Spending Trend subplot.
        """
        print("\n📊 Creating summary visualization (no monthly subplot)...")

        min_age, max_age = age_range
        sns.set_style("whitegrid")

        money_fmt = FuncFormatter(lambda x, pos: f"${x:,.0f}")

        fig = plt.figure(figsize=(20, 10))

        # 1. Total Spending by Purchase Type
        ax1 = plt.subplot(2, 3, 1)
        category_data = analysis['avg_by_category']['Total'].sort_values(ascending=True)
        category_data.plot(kind='barh', ax=ax1, color='steelblue')
        ax1.set_title('Total Spending by Purchase Type', fontsize=14, fontweight='bold')
        ax1.set_xlabel('Total Amount ($)')
        ax1.set_ylabel('Purchase Type')
        ax1.xaxis.set_major_formatter(money_fmt)

        # 2. Average Spending by Job Type (Top 10)
        ax2 = plt.subplot(2, 3, 2)
        job_data = analysis['avg_by_job']['Average'].head(10).sort_values(ascending=True)
        job_data.plot(kind='barh', ax=ax2, color='coral')
        ax2.set_title('Average Spending by Job Type (Top 10)', fontsize=14, fontweight='bold')
        ax2.set_xlabel('Average Amount ($)')
        ax2.set_ylabel('Job Type')
        ax2.xaxis.set_major_formatter(money_fmt)

        # 3. Average Spending by Education Level
        ax3 = plt.subplot(2, 3, 3)
        edu_data = analysis['avg_by_education']['Average'].sort_values(ascending=True)
        edu_data.plot(kind='barh', ax=ax3, color='purple')
        ax3.set_title('Average Spending by Education Level', fontsize=14, fontweight='bold')
        ax3.set_xlabel('Average Amount ($)')
        ax3.set_ylabel('Education Level')
        ax3.xaxis.set_major_formatter(money_fmt)

        # 4. Average Spending by Marital Status
        ax4 = plt.subplot(2, 3, 4)
        marital_data = analysis['avg_by_marital']['Average'].sort_values(ascending=True)
        marital_data.plot(kind='barh', ax=ax4, color='teal')
        ax4.set_title('Average Spending by Marital Status', fontsize=14, fontweight='bold')
        ax4.set_xlabel('Average Amount ($)')
        ax4.set_ylabel('Marital Status')
        ax4.xaxis.set_major_formatter(money_fmt)

        # 5. Transaction Distribution by Category (Pie)
        ax5 = plt.subplot(2, 3, 5)
        category_counts = analysis['avg_by_category']['Count']
        colors = plt.cm.Set3(range(len(category_counts)))
        ax5.pie(
            category_counts,
            labels=category_counts.index,
            autopct='%1.1f%%',
            startangle=90,
            colors=colors,
        )
        ax5.set_title('Transaction Distribution by Category', fontsize=14, fontweight='bold')

        plt.suptitle(
            f'Customer Spending Analysis: Ages {min_age}-{max_age}',
            fontsize=16,
            fontweight='bold',
            y=0.995,
        )
        plt.tight_layout()

        filename = f'{output_dir}spending_analysis_{min_age}_{max_age}.png'
        plt.savefig(filename, dpi=300, bbox_inches='tight')
        plt.close(fig)
        print(f"✓ Saved visualization: {filename}")

    def close(self):
        if self.connection:
            try:
                self.connection.close()
                print("\n✓ Connection closed")
            except Exception:
                pass


def main():
    """Interactive CLI entrypoint."""
    print("="*80)
    print("CUSTOMER SPENDING BEHAVIOR ANALYZER (NO MONTHLY SUBPLOT)")
    print("="*80)

    print("\nEnter the age range you want to analyze:")
    try:
        min_age = int(input("  Minimum age: "))
        max_age = int(input("  Maximum age: "))

        if min_age < 0 or max_age < 0 or min_age > max_age:
            print("✗ Invalid age range. Please ensure min_age <= max_age and both are positive.")
            return

        database = "bank_dw"
        print(f"\n📊 Using database: {database}")

    except ValueError:
        print("✗ Invalid input. Please enter numeric values for ages.")
        return

    analyzer = SpendingAnalyzer(host='localhost', port=10000, database=database)

    if not analyzer.connect():
        return

    try:
        df = analyzer.get_age_range_data(min_age, max_age)

        if df is None or len(df) == 0:
            print("✗ No data found for the specified age range.")
            analyzer.close()
            return

        analysis = analyzer.analyze_spending_behavior(df)

        # Print report once to screen
        analyzer.generate_text_report((min_age, max_age), analysis)

        # Save visuals
        analyzer.create_visualizations((min_age, max_age), analysis)

        # Save report to text file as well
        import sys
        from io import StringIO

        old_stdout = sys.stdout
        sys.stdout = StringIO()
        analyzer.generate_text_report((min_age, max_age), analysis)
        report_text = sys.stdout.getvalue()
        sys.stdout = old_stdout

        report_path = f"/home/hadoop/spending_report_{min_age}_{max_age}_{datetime.now().strftime('%Y%m%d_%H%M%S')}.txt"
        with open(report_path, "w") as f:
            f.write(report_text)
        print(f"\n✓ Report saved to: {report_path}")

    except Exception as e:
        print(f"\n✗ An error occurred: {e}")
        import traceback
        traceback.print_exc()
    finally:
        analyzer.close()
        print("\n" + "="*80)
        print("Analysis complete.")
        print("="*80)


if __name__ == "__main__":
    main()






